\scanpage{50}
\handwritten{\deleted{\uline{NB} En passant: pour que \unsure{deux|des} triades $\tau_1,\tau_2$ soient biliées, il f et s. qu'elles soient disjointes et appartiennent à un $=$ bitriangle, qui est \unsure{uniquemt|uniquement} déterminé, et est \unsure{aussi|aussitôt} l'unique}}

\handwritten{\textcircled{\scriptsize 11} $X$ est l'unique bitriangle contenant $\tau_i$ -- deux triades sont biliées ssi elles sont \inserted{disjointes et} contenues dans un $=$ bitriangle \deleted{(i.e. \unsure{l})}}

\handwritten{\textcircled{\scriptsize 12} Soit \deleted{\unsure{$X$ ou}} $(\tau_i)_{i\in I}$ un triangle de triades biliées, associé au bitriangle $X=\coprod_{i\in I}\tau_i\subset S$ avec sa coorientation $\{\tau_i\}\deleted{\unsure{\tau_2,\tau_3}}_{i\in I}$, \deleted{soit $(\bar\tau_i)_{i\in I}$ le triangle de} \inserted{$J$ l'ens des} \deleted{\unsure{triades}} \inserted{$(\tau'_j)_{j\in J}$} de $X$ de l'autre famille (de sorte que $X\xrightarrow{\sim}I\times J$). Soit $(\bar\tau_i)_{i\in I}$ le triangle de triades opposé, $\bar X=\bigcup_{i\in I}\bar\tau_i$, et $\bar J$ l'ens des triades de $\bar X$ de l'autre famille, donc $\bar X\xrightarrow{\sim}I\times\bar J$. Soit $\widetilde X=S-X-\bar X$, qui est de cardinal $9$ \inserted{$=27-2\cdot9$}. Alors}

\handwritten{a) $\widetilde X$ est un bitriangle. \deleted{\unsure{Il est|Soit}} L'application qui à tt hexagone de triades associe $\widetilde X$, est bijective. \deleted{\unsure{$S$}}}

% printed form
% Inverted relative to the manuscript. Long diagonal pen strokes cross the form.
\noindent\begin{minipage}[t]{.28\linewidth}
\deleted{Place}\\
\deleted{N\textsuperscript{o}}\par\smallskip

\deleted{UNIVERSITÉ DES SCIENCES ET TECHNIQUES DU LANGUEDOC}

\deleted{(1) -- Le candidat doit indiquer la nature de l'examen auquel il se présente.}

\deleted{(2) -- Le candidat inscrit ses nom, prénom, lieu et date de naissance en tête de la copie et signe lisiblement son nom à la fin de la composition.}

\end{minipage}\hfill\begin{minipage}[t]{.68\linewidth}
\deleted{ACADEMIE DE MONTPELLIER}\par\smallskip

\deleted{Examen de \textsuperscript{(1)}} \dotfill\par
\noindent\dotfill

\deleted{Nom et Prénoms du Candidat \textsuperscript{(2)} \dotfill}

\deleted{né à \dotfill{} département de \dotfill}

\deleted{le \dotfill{} 19}

\deleted{Montpellier, le \dotfill{} 19}

\deleted{COMPOSITION EN} \dotfill
\end{minipage}

% archive sheet number: 49
